\chapter{Aspectos éticos, seguridad y uso responsable}

\section{Enfoque general}

El uso responsable no se entiende en este trabajo como una advertencia añadida al final del desarrollo, sino como una condición que afecta a todo el recorrido del sistema. La consulta se acota antes de procesarse, el modelo trabaja mediante contratos definidos, el código se revisa antes de ejecutarse y la respuesta final conserva las limitaciones del análisis. Aun así, ninguna de estas medidas elimina por completo el riesgo de producir un cálculo incorrecto o una explicación inadecuada.

Este planteamiento coincide con los marcos actuales de gestión del riesgo en inteligencia artificial, que recomiendan considerar la fiabilidad, la transparencia y la supervisión durante todo el ciclo de vida, y no únicamente al evaluar la salida final~\cite{nist2024genai}. El Reglamento europeo de inteligencia artificial adopta también un enfoque basado en el riesgo y concede especial importancia a la supervisión humana, la robustez técnica y la información clara sobre las capacidades y limitaciones del sistema~\cite{euai2024regulation}. Estas referencias se utilizan como criterios de diseño; no se pretende establecer en este capítulo una clasificación jurídica del prototipo, que dependería de su finalidad y de las condiciones concretas de un posible despliegue.

\section{Finalidad prevista y usos excluidos}

El sistema está diseñado para responder preguntas sobre series financieras históricas ya disponibles. Puede calcular rentabilidad, volatilidad, drawdown, correlaciones, medias móviles, extremos u otras medidas descriptivas, y organizar los resultados en tablas o series estructuradas. Su finalidad es facilitar el análisis y mostrar de forma trazable cómo se ha llegado a una respuesta.

Quedan fuera de su alcance la predicción de precios, la recomendación de compra o venta, la gestión de carteras, la adaptación a la situación personal de un inversor y la ejecución de operaciones en el mercado. El sistema tampoco decide qué activo debe descargarse ni busca información por su cuenta. Trabaja con instrumentos identificados previamente y con rutas de datos proporcionadas en la entrada.

Esta delimitación es especialmente importante porque una respuesta bien redactada puede transmitir más seguridad de la que justifican los datos. Las métricas históricas describen lo ocurrido durante un periodo concreto, pero no garantizan que el comportamiento vaya a repetirse. Las advertencias dirigidas a inversores insisten precisamente en que las afirmaciones sobre resultados pasados deben presentarse con contexto y sin inducir expectativas injustificadas~\cite{investorgov2022performanceclaims}.

\begin{table}[H]
\centering
\small
\begin{tabularx}{\textwidth}{L{0.30\textwidth}Y}
\toprule
\textbf{Uso previsto} & \textbf{Uso no previsto} \\
\midrule
Análisis descriptivo de datos históricos seleccionados. & Predicción de precios o rentabilidades futuras. \\
Cálculo y explicación de métricas financieras. & Recomendaciones de inversión o asesoramiento personalizado. \\
Comparación trazable de activos y periodos. & Ejecución automática de órdenes o gestión de una cartera. \\
Apoyo a la revisión académica del flujo. & Sustitución de una validación financiera profesional. \\
\bottomrule
\end{tabularx}
\caption{Delimitación del uso responsable del sistema.}
\label{tab:uso_previsto_sistema}
\end{table}

\section{Interpretación financiera y supervisión humana}

El principal riesgo para el usuario no es únicamente que una cifra sea incorrecta. También puede ocurrir que un dato correcto se presente sin el contexto necesario, que se confunda una relación histórica con una explicación causal o que una respuesta fluida se acepte sin revisión. Esta tendencia a confiar en exceso en una salida automatizada resulta especialmente delicada cuando el contenido se refiere a decisiones económicas.

El diseño intenta reducir este riesgo separando planificación, cálculo e interpretación. El agente final no recibe libertad para reconstruir el análisis: trabaja con el plan y con la salida obtenida al ejecutar el código. Además, se le pide distinguir entre datos observados, lectura interpretativa y limitaciones, y se prohíbe incorporar información externa o completar cifras ausentes. Esta separación facilita comprobar si la redacción se corresponde con los resultados.

No obstante, los tres agentes utilizan el mismo modelo. Por ello, la división funcional mejora la trazabilidad, pero no equivale a contar con tres revisores independientes. Un supuesto equivocado puede propagarse desde el plan hasta el código y conservar una apariencia coherente en la respuesta final. La revisión humana sigue siendo necesaria para confirmar la identidad del instrumento, la adecuación de las métricas, las unidades, el periodo analizado y la interpretación de los resultados.

La supervisión debe ser proporcional al uso. En el contexto académico resulta suficiente conservar los artefactos y revisar los casos relevantes. En una aplicación dirigida a usuarios externos sería necesario mostrar con claridad que la respuesta ha sido generada con ayuda de inteligencia artificial, facilitar el acceso a sus fuentes y limitaciones, y permitir que una persona descarte el resultado cuando detecte una anomalía.

\section{Calidad, procedencia y alcance de los datos}

La calidad de la respuesta está condicionada por la calidad de los datos de entrada. La batería utiliza series descargadas previamente y almacenadas de forma local. Esta decisión favorece la reproducibilidad y evita que el sistema elija automáticamente un instrumento distinto del solicitado, pero también limita su autonomía y hace necesario preparar las fuentes antes de formular la consulta.

La identificación del activo no siempre es trivial. Una expresión como ``oro'' puede aludir a un precio de referencia, a un contrato de futuros o a un fondo cotizado. Del mismo modo, dos series sobre un mismo mercado pueden diferir por moneda, horario, ajustes o calendario. El sistema no resuelve por sí solo estas ambigüedades; presupone que el conjunto de datos asociado a la consulta ha sido seleccionado correctamente.

Los datos proceden de descargas realizadas mediante \texttt{yfinance}, una herramienta abierta que no está afiliada ni respaldada por Yahoo~\cite{yfinance,yahoofinancehistorical}. El prototipo tampoco incorpora noticias, información fundamental, variables macroeconómicas, costes de transacción ni datos en tiempo real. En comparaciones entre mercados diferentes no siempre se homogeneizan monedas y calendarios. Estas ausencias no invalidan un análisis descriptivo, pero sí restringen las conclusiones que pueden extraerse.

La batería define los componentes esperados de cada respuesta, aunque no contiene una referencia numérica validada manualmente para todas las métricas. En consecuencia, completar el flujo demuestra que el sistema ha podido planificar, generar, ejecutar e interpretar una solución, pero no certifica por sí solo la corrección financiera de cada cifra. Para un uso de mayor impacto sería necesario contrastar los cálculos con implementaciones de referencia y añadir controles específicos para cada familia de métricas.

\section{Seguridad del código generado}

El código producido por un LLM debe considerarse una entrada no fiable. OWASP identifica como riesgos relevantes tanto el tratamiento insuficiente de las salidas del modelo como la concesión de más funciones, permisos o autonomía de los necesarios~\cite{owasp2025improperoutput,owasp2025excessiveagency}. En este sistema, el modelo puede proponer un programa, pero no decide directamente si debe ejecutarse.

La validación analiza el árbol sintáctico del código y exige una función principal y una salida JSON. Se limita el conjunto de importaciones y se rechazan llamadas como \texttt{eval}, \texttt{exec}, \texttt{compile}, \texttt{open} o \texttt{\_\_import\_\_}. También se bloquean módulos asociados a procesos, red y manipulación libre del sistema de archivos. Los prompts refuerzan estas condiciones al exigir las funciones de carga previstas y las rutas recibidas en la entrada.

Si el código supera esta revisión, se ejecuta en un proceso independiente con un límite de 60 segundos. Se conservan el programa, el payload, la salida estándar, la salida de error y el código de retorno. Cuando aparece un rechazo o un fallo de ejecución, puede realizarse un único intento de reparación; la nueva versión vuelve a pasar por la misma validación antes de ejecutarse. Si no supera los controles, el flujo termina con un error explícito y no fabrica una respuesta financiera alternativa.

Estas medidas reducen la superficie de riesgo, pero no constituyen una sandbox completa. El proceso comparte el entorno de trabajo, hereda variables del proceso principal y no dispone de límites específicos de memoria o CPU. La revisión estática tampoco demuestra que el cálculo sea correcto ni garantiza que todas las formas posibles de acceso a recursos queden cubiertas. Por este motivo, el prototipo no debe ejecutar código generado a partir de entradas no confiables en un entorno expuesto.

\begin{table}[H]
\centering
\small
\begin{tabularx}{\textwidth}{L{0.27\textwidth}L{0.34\textwidth}Y}
\toprule
\textbf{Riesgo} & \textbf{Control actual} & \textbf{Riesgo residual} \\
\midrule
Código sintácticamente inválido & Análisis previo y reparación limitada. & La reparación también puede fallar. \\
Uso de operaciones peligrosas & Lista de importaciones y llamadas bloqueadas. & La revisión estática no equivale a aislamiento. \\
Ejecución indefinida & Proceso separado y límite temporal. & No hay cuotas específicas de memoria o CPU. \\
Salida no interpretable & Contrato JSON y registro de \texttt{stdout} y \texttt{stderr}. & Un JSON válido puede contener cálculos equivocados. \\
Respuesta no respaldada & Interpretación basada en el plan y la salida ejecutada. & El modelo puede omitir contexto o expresar mal una cifra. \\
\bottomrule
\end{tabularx}
\caption{Controles de seguridad y riesgos que permanecen abiertos.}
\label{tab:controles_riesgos_residuales}
\end{table}

\section{Trazabilidad y rendición de cuentas}

Cada ejecución conserva la entrada, el plan, el código, el resultado, los avisos, los errores y la respuesta final. Esta información permite reconstruir el recorrido y atribuir un fallo a la etapa en la que se produjo. La trazabilidad no convierte automáticamente el sistema en fiable, pero evita que la única evidencia sea un texto final convincente.

La batería añade una preevaluación heurística de cinco dimensiones para facilitar la revisión. Estas puntuaciones ayudan a localizar casos problemáticos, aunque no sustituyen el criterio humano. Algunas dimensiones se aproximan mediante la presencia de términos, la longitud de la respuesta o el estado del flujo, de modo que una puntuación alta no demuestra que cada cálculo sea correcto.

En un despliegue real sería necesario definir quién responde de la selección de datos, quién aprueba las métricas, quién revisa las incidencias y durante cuánto tiempo se conservan los registros. También deberían versionarse el modelo, los prompts, las dependencias y el catálogo de datos para poder reproducir una respuesta concreta.

\section{Privacidad, confidencialidad y credenciales}

Los conjuntos utilizados contienen precios y volúmenes históricos, no información personal. Las claves de acceso al servicio LLM se almacenan en el fichero local \texttt{.env}, excluido del control de versiones. Los directorios que contienen código generado, payloads, registros e informes de evaluación también se excluyen del repositorio.

Sin embargo, las llamadas al modelo transmiten la consulta, información estructurada sobre los activos, las rutas de los datos, el plan y, en la fase final, los resultados calculados. Además, los registros locales conservan una parte importante de ese contenido. En la batería actual esta información no es sensible, pero podría serlo si una versión futura aceptara carteras, estrategias internas o datos identificables de usuarios.

En ese escenario sería necesario aplicar minimización de datos, sustituir rutas locales por identificadores neutros, establecer plazos de conservación, controlar el acceso a los registros y documentar las condiciones de tratamiento del proveedor. Si se incorporasen datos personales, estos requisitos deberían alinearse con los principios de limitación de finalidad, minimización, exactitud e integridad establecidos por el Reglamento General de Protección de Datos~\cite{gdpr2016}.

\section{Coste computacional y sostenibilidad}

El trabajo no entrena ni ajusta un modelo propio. Reutiliza un modelo ya desplegado y reserva las llamadas generativas para tres funciones concretas: planificación, generación de código e interpretación. Los cálculos financieros se realizan localmente sobre datos almacenados.

Esto no significa que el coste ambiental sea despreciable. Cada caso requiere normalmente varias inferencias y puede añadir una llamada de reparación. La infraestructura que sirve el modelo consume recursos que no se han medido en este estudio, por lo que no sería riguroso afirmar una reducción cuantificada de energía o emisiones.

Las medidas adoptadas son principalmente operativas: temperatura baja, límites de salida y tiempo, compactación del contexto enviado al agente final, reparación acotada y posibilidad de continuar una batería sin repetir casos ya completados. Estas decisiones contienen el consumo innecesario y el coste económico, aunque una evaluación ambiental sólida exigiría medir energía, hardware, duración y carga del servicio.

\section{Condiciones para un uso futuro}

El prototipo es adecuado como demostrador académico y como entorno de evaluación de un flujo LLM trazable. No está preparado para actuar como herramienta profesional autónoma. Antes de ofrecerlo a usuarios externos sería necesario, al menos, incorporar aislamiento real para el código, permisos mínimos, límites de recursos, gestión formal de registros y credenciales, validaciones financieras de referencia y revisión humana en los casos de mayor impacto.

También debería mantenerse una separación clara entre análisis y acción. Conectar el sistema a una cuenta de inversión, permitirle enviar órdenes o personalizar decisiones a partir de datos económicos de una persona cambiaría sustancialmente el nivel de riesgo y exigiría una evaluación jurídica, técnica y organizativa específica.

La conclusión práctica es deliberadamente prudente: el sistema puede ayudar a organizar y explicar un análisis histórico, pero la responsabilidad de comprobar los datos, interpretar el alcance y decidir si el resultado es utilizable continúa siendo humana.
